A patch can attract strong attention while returning a short value vector. We use the fast-weight view of attention to test how value lengths affect spatial maps and how those maps relate to classification. Norm weighting improves object localisation and human alignment in supervised DeiT. On Oxford-IIIT Pet, assigning each image's value lengths to positions ranked by another image reduces localisation. For classification, we directly remove individual patch contributions to the final classification token. Restricting the readout to this query improves agreement with changes in the class margin and reduces localisation on the same images. The spatial assignment of value lengths and the choice of receiving query thus explain different aspects of an attention map.
